\section{Available output can still create losses and obligations}
\label{sec:apple-quality-inventory}

Obtaining components and completing delivery do not establish that a product
can safely serve its intended use. Apple's defects discussion covers its
hardware and software, third-party components and products, and service
outages, slowdowns and errors. It reports that services have sometimes failed
to perform as anticipated. It also warns that defects can create personal-
injury, property or environmental risks, that specialized uses such as health
can increase those risks, and that complex technologies such as artificial
intelligence can expose users to harmful or inaccurate content. Detection and
repair of every defect are not assured \citep[Risk Factors, printed p.~9]{apple2024}.
For the model, a delivered unit, an acceptable unit and a repaired unit are
different states. The source identifies possible transitions between them,
not a defect rate or a particular incident.

The financial consequences can extend well beyond inventory. Apple names
product liability, recalls, replacement or modification, write-offs of
inventory, property, plant and equipment or intangible assets, warranty and
other expenses, litigation costs and regulatory fines. Quality can also affect
reputation, acceptance, demand and the timing of new offerings. These effects
must not all become one inventory expense. A customer obligation and a possible
supplier recovery are also separate. The outsourcing passage says partner
arrangements may include defect-expense reimbursement while Apple generally
remains responsible to the consumer for warranty and out-of-warranty service
\citep[Risk Factors, printed p.~8]{apple2024}. Identifying a supplier does not
establish a reimbursement right, amount or receipt date.

Excess stock raises a different question even when its quality is adequate.
The inventory-risk passage says Apple records write-downs for product and
component inventories that are obsolete, exceed anticipated demand, or have
cost above net realizable value. It separately describes necessary cancellation
fee reserves for orders of excess products and components. Inventory already
held and an order not yet received are different exposures; cancellation
depends on the actual order terms. Note 1's first-in, first-out policy identifies
the reported cost-flow method, rather than a lot's realizable value or a
cancellation charge \citep[Risk Factors, printed p.~9; Note 1, printed p.~34]{apple2024}.

Apple also describes impairment reviews when circumstances indicate that
long-lived assets, including capital assets held at suppliers and inventory
prepayments, may not be recoverable. It reports a write-down for excess carrying
value over fair value if impairment is determined. This is the issuer's
description of its assessment; an executable calculation still needs the
actual asset classification, assessment basis and applicable accounting method.
The company's belief that the named assets and commitments are currently
recoverable is accompanied by no assurance against future charges. It neither
proves a future impairment nor permits the model to erase that possibility.
An inventory cost-flow calculation, a cancellation obligation and an asset
impairment assessment cannot be substituted for one another.

The link to product transitions is the forecast used before delivery. Apple
orders components and builds inventory ahead of announcements and shipments;
manufacturing purchase obligations cover forecast requirements, typically for
periods up to 150 days. It warns of excess or insufficient output, or failure
to use firm commitments fully, when forecasts are wrong. The qualifier
``typically'' prevents that period from becoming a universal contract limit.
Nor does the description make the separately reported Note 12 obligations
with remaining terms above one year the same population as manufacturing
obligations \citep[Risk Factors, printed p.~9; Note 12, printed p.~45]{apple2024}.

Consider an explicitly assumed reduction in expected demand. Two companies
could hold the same reported inventory amount while one holds saleable
components for the new offering and the other holds incompatible earlier
components, defective finished goods or commitments it cannot cancel cheaply.
Their feasible responses and possible losses differ. Lot identity and cost,
quality, compatibility, forecast dates, anticipated demand, realizable value,
order terms and customer or supplier obligations distinguish the histories.
Recognition and measurement then require their own evidence. The disclosure
does not supply those facts or probabilities, so the model preserves the
conditional histories without inventing a reserve or treating the forecast as
observed customer demand.

\subsection{Service provision requires continuing access}
\label{sec:apple-service-access}

Adequate components and saleable inventory still leave a service question:
can the company continue offering the software and content its customers want?
Apple says its products and services require third-party intellectual property,
while giving no assurance that the necessary licences will be available on
commercially reasonable terms or at all. Its digital-content arrangements can
be short and need not renew on similar terms. Content providers can also require
new rights-management or security technology before access can continue
\citep[Risk Factors, printed pp.~9--10]{apple2024}. The relevant resource is
therefore partly a dated right to offer particular content or functionality.
Physical capacity alone cannot establish that right.

Developer support has a separate economic cause. Apple believes customers'
hardware purchases depend partly on available applications, whose continued
development and maintenance are not assured. It describes developers comparing
platforms by market position, anticipated revenue, expected growth and the cost
of developing and updating software \citep[Risk Factors, printed p.~10]{apple2024}.
An application can consequently become unavailable or unsuitable even while
the hardware remains usable. A scenario needs the affected applications and
developers' response; an estimate of component capacity cannot supply either.

Access also depends on how a service uses data and takes payment. Apple reports
restrictions on collecting, using and transferring personal information,
including transfers within its international group. It separately warns that
failure to follow payment-card security standards could increase transaction
costs or remove the ability to process cards
\citep[Risk Factors, printed p.~15]{apple2024}. These issuer descriptions
identify possible constraints. Which constraint applies depends on the
data, entity, jurisdiction, agreement and effective rule.

Suppose a content agreement permits a specified offering only until its expiry
date. Under one explicitly assumed history it renews; under another it expires
without replacement. Identical physical resources then permit different
offerings. The comparison needs the service, customer population, territory,
contract term and allowed substitutions, together with any new cost and demand
response. A financial calculation must apply the relevant revenue, payment
and remedy terms to the service actually provided. Apple's risk statements
supply neither those contract facts nor probabilities for the alternatives.
The historical record must preserve the continuing relationships as well as
the reported balances if it is to support these different futures.
